% Options for packages loaded elsewhere
\PassOptionsToPackage{unicode}{hyperref}
\PassOptionsToPackage{hyphens}{url}
\documentclass[
]{article}
\usepackage{xcolor}
\usepackage[margin=1in]{geometry}
\usepackage{amsmath,amssymb}
\setcounter{secnumdepth}{-\maxdimen} % remove section numbering
\usepackage{iftex}
\ifPDFTeX
  \usepackage[T1]{fontenc}
  \usepackage[utf8]{inputenc}
  \usepackage{textcomp} % provide euro and other symbols
\else % if luatex or xetex
  \usepackage{unicode-math} % this also loads fontspec
  \defaultfontfeatures{Scale=MatchLowercase}
  \defaultfontfeatures[\rmfamily]{Ligatures=TeX,Scale=1}
\fi
\usepackage{lmodern}
\ifPDFTeX\else
  % xetex/luatex font selection
\fi
% Use upquote if available, for straight quotes in verbatim environments
\IfFileExists{upquote.sty}{\usepackage{upquote}}{}
\IfFileExists{microtype.sty}{% use microtype if available
  \usepackage[]{microtype}
  \UseMicrotypeSet[protrusion]{basicmath} % disable protrusion for tt fonts
}{}
\makeatletter
\@ifundefined{KOMAClassName}{% if non-KOMA class
  \IfFileExists{parskip.sty}{%
    \usepackage{parskip}
  }{% else
    \setlength{\parindent}{0pt}
    \setlength{\parskip}{6pt plus 2pt minus 1pt}}
}{% if KOMA class
  \KOMAoptions{parskip=half}}
\makeatother
\usepackage{color}
\usepackage{fancyvrb}
\newcommand{\VerbBar}{|}
\newcommand{\VERB}{\Verb[commandchars=\\\{\}]}
\DefineVerbatimEnvironment{Highlighting}{Verbatim}{commandchars=\\\{\}}
% Add ',fontsize=\small' for more characters per line
\usepackage{framed}
\definecolor{shadecolor}{RGB}{248,248,248}
\newenvironment{Shaded}{\begin{snugshade}}{\end{snugshade}}
\newcommand{\AlertTok}[1]{\textcolor[rgb]{0.94,0.16,0.16}{#1}}
\newcommand{\AnnotationTok}[1]{\textcolor[rgb]{0.56,0.35,0.01}{\textbf{\textit{#1}}}}
\newcommand{\AttributeTok}[1]{\textcolor[rgb]{0.13,0.29,0.53}{#1}}
\newcommand{\BaseNTok}[1]{\textcolor[rgb]{0.00,0.00,0.81}{#1}}
\newcommand{\BuiltInTok}[1]{#1}
\newcommand{\CharTok}[1]{\textcolor[rgb]{0.31,0.60,0.02}{#1}}
\newcommand{\CommentTok}[1]{\textcolor[rgb]{0.56,0.35,0.01}{\textit{#1}}}
\newcommand{\CommentVarTok}[1]{\textcolor[rgb]{0.56,0.35,0.01}{\textbf{\textit{#1}}}}
\newcommand{\ConstantTok}[1]{\textcolor[rgb]{0.56,0.35,0.01}{#1}}
\newcommand{\ControlFlowTok}[1]{\textcolor[rgb]{0.13,0.29,0.53}{\textbf{#1}}}
\newcommand{\DataTypeTok}[1]{\textcolor[rgb]{0.13,0.29,0.53}{#1}}
\newcommand{\DecValTok}[1]{\textcolor[rgb]{0.00,0.00,0.81}{#1}}
\newcommand{\DocumentationTok}[1]{\textcolor[rgb]{0.56,0.35,0.01}{\textbf{\textit{#1}}}}
\newcommand{\ErrorTok}[1]{\textcolor[rgb]{0.64,0.00,0.00}{\textbf{#1}}}
\newcommand{\ExtensionTok}[1]{#1}
\newcommand{\FloatTok}[1]{\textcolor[rgb]{0.00,0.00,0.81}{#1}}
\newcommand{\FunctionTok}[1]{\textcolor[rgb]{0.13,0.29,0.53}{\textbf{#1}}}
\newcommand{\ImportTok}[1]{#1}
\newcommand{\InformationTok}[1]{\textcolor[rgb]{0.56,0.35,0.01}{\textbf{\textit{#1}}}}
\newcommand{\KeywordTok}[1]{\textcolor[rgb]{0.13,0.29,0.53}{\textbf{#1}}}
\newcommand{\NormalTok}[1]{#1}
\newcommand{\OperatorTok}[1]{\textcolor[rgb]{0.81,0.36,0.00}{\textbf{#1}}}
\newcommand{\OtherTok}[1]{\textcolor[rgb]{0.56,0.35,0.01}{#1}}
\newcommand{\PreprocessorTok}[1]{\textcolor[rgb]{0.56,0.35,0.01}{\textit{#1}}}
\newcommand{\RegionMarkerTok}[1]{#1}
\newcommand{\SpecialCharTok}[1]{\textcolor[rgb]{0.81,0.36,0.00}{\textbf{#1}}}
\newcommand{\SpecialStringTok}[1]{\textcolor[rgb]{0.31,0.60,0.02}{#1}}
\newcommand{\StringTok}[1]{\textcolor[rgb]{0.31,0.60,0.02}{#1}}
\newcommand{\VariableTok}[1]{\textcolor[rgb]{0.00,0.00,0.00}{#1}}
\newcommand{\VerbatimStringTok}[1]{\textcolor[rgb]{0.31,0.60,0.02}{#1}}
\newcommand{\WarningTok}[1]{\textcolor[rgb]{0.56,0.35,0.01}{\textbf{\textit{#1}}}}
\usepackage{graphicx}
\makeatletter
\newsavebox\pandoc@box
\newcommand*\pandocbounded[1]{% scales image to fit in text height/width
  \sbox\pandoc@box{#1}%
  \Gscale@div\@tempa{\textheight}{\dimexpr\ht\pandoc@box+\dp\pandoc@box\relax}%
  \Gscale@div\@tempb{\linewidth}{\wd\pandoc@box}%
  \ifdim\@tempb\p@<\@tempa\p@\let\@tempa\@tempb\fi% select the smaller of both
  \ifdim\@tempa\p@<\p@\scalebox{\@tempa}{\usebox\pandoc@box}%
  \else\usebox{\pandoc@box}%
  \fi%
}
% Set default figure placement to htbp
\def\fps@figure{htbp}
\makeatother
\setlength{\emergencystretch}{3em} % prevent overfull lines
\providecommand{\tightlist}{%
  \setlength{\itemsep}{0pt}\setlength{\parskip}{0pt}}
\usepackage{bookmark}
\IfFileExists{xurl.sty}{\usepackage{xurl}}{} % add URL line breaks if available
\urlstyle{same}
\hypersetup{
  pdftitle={Logistic Regression},
  hidelinks,
  pdfcreator={LaTeX via pandoc}}

\title{Logistic Regression}
\author{}
\date{\vspace{-2.5em}2026-04-17}

\begin{document}
\maketitle

\subsection{Introduction}\label{introduction}

Logistic regression models binary outcomes via the logit link, mapping
the probability of ``event'' to the entire real line so that linear
regression machinery applies on the log-odds scale. It is the most
widely used GLM in clinical, epidemiological, and behavioural research,
producing odds ratios that directly answer the question ``how does the
odds of the outcome change with this covariate?'' The fit is by
iteratively reweighted least squares (IRLS), maximising the binomial
likelihood.

\subsection{Prerequisites}\label{prerequisites}

A working understanding of the binomial distribution, maximum likelihood
estimation, the logit link, and the distinction between odds ratios and
relative risks.

\subsection{Theory}\label{theory}

The model is

\[Y \sim \mathrm{Bernoulli}(p), \qquad \mathrm{logit}(p) = \log\frac{p}{1-p} = \beta_0 + \sum_j \beta_j X_j.\]

IRLS produces \(\hat\beta\) on the log-odds scale; exponentiating gives
odds ratios, which are interpretable as the multiplicative change in the
odds of the outcome per unit increase in the covariate, all else equal.
For rare events (prevalence below \textasciitilde10 \%), the odds ratio
approximates the relative risk; for common events the two diverge and
odds ratios overstate the relative risk.

\subsection{Assumptions}\label{assumptions}

Binary outcome, independent observations, linearity on the logit scale
(the log-odds is a linear function of the covariates), absence of severe
multicollinearity, and adequate events per variable (EPV \(\ge 10\) as a
working rule for stable estimation).

\subsection{R Implementation}\label{r-implementation}

\begin{Shaded}
\begin{Highlighting}[]
\FunctionTok{library}\NormalTok{(broom); }\FunctionTok{library}\NormalTok{(pROC); }\FunctionTok{library}\NormalTok{(performance)}

\NormalTok{d }\OtherTok{\textless{}{-}}\NormalTok{ mtcars}
\NormalTok{d}\SpecialCharTok{$}\NormalTok{am }\OtherTok{\textless{}{-}} \FunctionTok{factor}\NormalTok{(d}\SpecialCharTok{$}\NormalTok{am)}

\NormalTok{fit }\OtherTok{\textless{}{-}} \FunctionTok{glm}\NormalTok{(am }\SpecialCharTok{\textasciitilde{}}\NormalTok{ wt }\SpecialCharTok{+}\NormalTok{ hp, }\AttributeTok{data =}\NormalTok{ d, }\AttributeTok{family =}\NormalTok{ binomial)}
\FunctionTok{tidy}\NormalTok{(fit, }\AttributeTok{conf.int =} \ConstantTok{TRUE}\NormalTok{, }\AttributeTok{exponentiate =} \ConstantTok{TRUE}\NormalTok{)}
\FunctionTok{glance}\NormalTok{(fit)}

\CommentTok{\# Discrimination}
\NormalTok{pred }\OtherTok{\textless{}{-}} \FunctionTok{predict}\NormalTok{(fit, }\AttributeTok{type =} \StringTok{"response"}\NormalTok{)}
\FunctionTok{roc}\NormalTok{(d}\SpecialCharTok{$}\NormalTok{am, pred, }\AttributeTok{quiet =} \ConstantTok{TRUE}\NormalTok{) }\SpecialCharTok{|\textgreater{}} \FunctionTok{auc}\NormalTok{()}

\FunctionTok{check\_model}\NormalTok{(fit)}
\end{Highlighting}
\end{Shaded}

\subsection{Output \& Results}\label{output-results}

\texttt{tidy(fit,\ exponentiate\ =\ TRUE)} returns odds ratios with
confidence intervals; \texttt{glance()} provides AIC, deviance, and
pseudo-\(R^2\) summaries. The AUC quantifies discrimination on the same
data; for honest assessment, cross-validate or hold out a test set.

\subsection{Interpretation}\label{interpretation}

A reporting sentence: ``A multivariable logistic regression of
transmission type on weight and horsepower estimated that each
additional 1000 lbs of weight reduced the odds of manual transmission by
a factor of 0.10 (95 \% CI 0.02 to 0.43, \(p = 0.003\)), adjusted for
horsepower; the model's in-sample AUC was 0.93 (95 \% CI 0.84 to
1.00).'' Always report odds ratios, not raw coefficients on the log-odds
scale.

\subsection{Practical Tips}\label{practical-tips}

\begin{itemize}
\tightlist
\item
  Report odds ratios with confidence intervals; coefficients on the
  log-odds scale are rarely the most communicative summary for clinical
  readers.
\item
  For rare events (prevalence \textless{} 10 \%), OR ≈ RR; for common
  events, OR overstates RR. Switch to \texttt{logbin} or modified
  Poisson for direct relative-risk regression when needed.
\item
  Small samples can produce complete separation, where MLE drifts to
  \(\pm\infty\); use Firth correction (\texttt{logistf::logistf}) or
  weakly informative Bayesian priors (\texttt{brms}) for finite
  estimates.
\item
  Calibration matters as much as discrimination; assess with
  \texttt{performance::check\_model()} calibration plots and the
  Hosmer-Lemeshow test (with caveats about its known limitations).
\item
  For matched case-control data, use conditional logistic regression
  (\texttt{survival::clogit}) which preserves the matched-pair
  structure.
\item
  For high-dimensional or correlated predictors, switch to penalised
  logistic regression via \texttt{glmnet(family\ =\ "binomial")}.
\end{itemize}

\subsection{R Packages Used}\label{r-packages-used}

Base R \texttt{glm()} with \texttt{family\ =\ binomial};
\texttt{broom::tidy()} and \texttt{glance()} for tidy summaries;
\texttt{pROC} for ROC and AUC; \texttt{performance::check\_model()} for
diagnostic plots; \texttt{logistf} for Firth-penalised logistic
regression; \texttt{glmnet} for elastic-net logistic regression;
\texttt{survival::clogit} for matched designs.

\subsection{For Reviewers}\label{for-reviewers}

\textbf{What to look for in a paper using this method.}

\begin{itemize}
\tightlist
\item
  \textbf{Common misapplications.}
\item
  \textbf{Diagnostics that should be reported but often aren't.}
\item
  \textbf{Red flags in tables and figures.}
\item
  \textbf{What to verify.}
\item
  \textbf{What an adequate Methods paragraph must contain.}
\end{itemize}

\subsection{See also --- labs in this
chapter}\label{see-also-labs-in-this-chapter}

\begin{itemize}
\tightlist
\item
  \href{labs/lab-correlation-vs-regression-model-i-ii-regression.Rmd}{Correlation
  vs regression; Model-I/II regression}
\item
  \href{labs/lab-simple-linear-regression.Rmd}{Simple linear regression}
\item
  \href{labs/lab-multiple-regression-confounding-interaction-centring.Rmd}{Multiple
  regression: confounding, interaction, centring}
\item
  \href{labs/lab-diagnostics-residuals-qq-leverage-cook-s-distance-vif.Rmd}{Diagnostics:
  residuals, QQ, leverage, Cook's distance, VIF}
\item
  \href{labs/lab-robust-regression-weighted-ls-hc-sandwich-standard-errors.Rmd}{Robust
  regression, weighted LS, HC (sandwich) standard errors}
\item
  \href{labs/lab-one-way-anova-and-contrasts-emmeans.Rmd}{One-way ANOVA
  and contrasts (emmeans)}
\item
  \href{labs/lab-two-way-factorial-anova-with-interaction.Rmd}{Two-way /
  factorial ANOVA with interaction}
\item
  \href{labs/lab-rcbd-and-repeated-measures-mixed-models.Rmd}{RCBD and
  repeated measures → mixed models}
\item
  \href{labs/lab-gams-with-mgcv.Rmd}{GAMs with mgcv}
\item
  \href{labs/lab-non-linear-regression-with-nls.Rmd}{Non-linear
  regression with nls}
\item
  \href{labs/lab-logistic-regression-binomial-glm.Rmd}{Logistic
  regression (binomial GLM)}
\item
  \href{labs/lab-ancova-in-rcts-baseline-adjustment.Rmd}{ANCOVA in RCTs
  (baseline adjustment)}
\item
  \href{labs/lab-ordinal-and-multinomial-regression.Rmd}{Ordinal and
  multinomial regression}
\item
  \href{labs/lab-poisson-and-negative-binomial-regression.Rmd}{Poisson
  and negative-binomial regression}
\item
  \href{labs/lab-calibration-discrimination-roc-auc-brier-score.Rmd}{Calibration,
  discrimination, ROC/AUC, Brier score}
\item
  \href{labs/lab-dichotomisation-change-scores-regression-to-the-mean.Rmd}{Dichotomisation,
  change scores, regression to the mean}
\item
  \href{labs/lab-decision-curves-nri-idi.Rmd}{Decision curves, NRI, IDI}
\end{itemize}

\end{document}
